\documentclass[11pt,a4paper]{article}
\usepackage[utf8]{inputenc}
\usepackage[T1]{fontenc}
\usepackage[catalan]{babel}
\usepackage{lmodern}
\usepackage[sfdefault,lining]{FiraSans}
\usepackage[a4paper,top=2.3cm,bottom=2.2cm,left=1.8cm,right=1.8cm,headheight=1.3cm,headsep=0.4cm,footskip=1cm]{geometry}
\usepackage{amsmath,amssymb}
\usepackage{icomma}
\usepackage{graphicx}
\usepackage[table]{xcolor}
\usepackage{tikz}
\usepackage{fontawesome5}
\usepackage{ccicons}
\usepackage{fancyhdr}
\usepackage{lastpage}
\usepackage{tasks}
\usepackage{needspace}
\usepackage[most]{tcolorbox}
\usepackage{array}
\usepackage[hidelinks,pdftitle={Fitxa d'exercicis de la unitat 1: Nombres naturals, enters i fraccionaris. 4t d'ESO Matemàtiques A},pdfsubject={Matemàtiques A 4t ESO},pdfkeywords={nombres naturals, operacions combinades, potències, arrels, nombres enters, valor absolut, fraccions, fraccions equivalents, simplificació, amplificació, problemes amb fraccions, propietats de les potències}]{hyperref}

\definecolor{navy}{HTML}{1B2748}
\definecolor{accent}{HTML}{E8603A}
\definecolor{teal}{HTML}{1F9E90}
\definecolor{soft}{HTML}{EEF1F8}
\definecolor{muted}{HTML}{5D6685}

\setlength{\parindent}{0pt}
\setlength{\parskip}{3pt}
\settasks{label=\alph*), label-format=\bfseries\color{navy}, label-width=1.3em, item-indent=1.8em, column-sep=0.6em, after-item-skip=2pt, before-skip=3pt, after-skip=4pt}

% ---------- capçalera i peu ----------
\newcommand{\llicencia}{\href{https://creativecommons.org/licenses/by-nc-sa/4.0/deed.ca}{\ccbyncsa}\ \ Llicència Creative Commons Reconeixement-NoComercial-CompartirIgual 4.0 Internacional (CC BY-NC-SA 4.0)}
\pagestyle{fancy}
\fancyhf{}
\fancyhead[L]{\raisebox{-0.35\height}{\includegraphics[height=0.95cm]{logo-gva.png}}}
\fancyhead[C]{\small\color{muted}Fitxa d'exercicis · Unitat 1: Nombres naturals, enters i fraccionaris · 4t ESO A}
\fancyhead[R]{\raisebox{-0.35\height}{\includegraphics[height=0.95cm]{logo-vermellar.png}}}
\fancyfoot[L]{\scriptsize\color{muted}\llicencia}
\fancyfoot[R]{\small\color{muted}\thepage}
\renewcommand{\headrulewidth}{0.4pt}
\renewcommand{\headrule}{\color{muted!40}\hrule height \headrulewidth}
\fancypagestyle{primera}{\fancyhf{}\renewcommand{\headrulewidth}{0pt}
  \fancyfoot[L]{\scriptsize\color{muted}\llicencia}
  \fancyfoot[R]{\small\color{muted}\thepage}}

% ---------- nivells i exercicis ----------
\newcommand{\estrelles}[1]{{\color{accent}\ifcase#1\or\faStar\faStar[regular]\faStar[regular]\or\faStar\faStar\faStar[regular]\or\faStar\faStar\faStar\fi}}
\newcounter{exer}
\newcommand{\ex}[1]{\par\needspace{4\baselineskip}\medskip\stepcounter{exer}\noindent{\color{navy}\bfseries\theexer.}\ \ {\small\estrelles{#1}}\ \ }

\newcommand{\seccio}[2]{\par\needspace{10\baselineskip}\bigskip
  \begin{tcolorbox}[colback=navy,colframe=navy,coltext=white,arc=2mm,boxrule=0pt,left=3mm,right=3mm,top=1.2mm,bottom=1.2mm]
  {\large\bfseries #1\quad #2}
  \end{tcolorbox}\vspace{-2pt}}

\newcolumntype{C}[1]{>{\centering\arraybackslash}p{#1}}

% quadre de recordatori (només a la fitxa d'exercicis)
\newcommand{\recorda}[1]{\par\smallskip\begin{tcolorbox}[colback=accent!7,colframe=accent!55,arc=2mm,boxrule=0.5pt,left=2.5mm,right=2.5mm,top=1mm,bottom=1mm]\small{\color{accent}\faLightbulb[regular]}\ \textbf{Recorda:}\ #1\end{tcolorbox}}
\newcommand{\res}[1]{#1}

\def\fmneed{11}% que el títol d'una secció no quede sol al final del full
% =====================================================================
% Capçalera FentMates de les fitxes: un color per a cada curs.
% Els colors són els de la pàgina inicial (LEVELS a assets/js/materials.js);
% el logo ve de _eines/logo/pdf/ (python3 _eines/logo/genera.py).
%
% Ús, al final del preàmbul de la unitat:
%   % =====================================================================
% Capçalera FentMates de les fitxes: un color per a cada curs.
% Els colors són els de la pàgina inicial (LEVELS a assets/js/materials.js);
% el logo ve de _eines/logo/pdf/ (python3 _eines/logo/genera.py).
%
% Ús, al final del preàmbul de la unitat:
%   % =====================================================================
% Capçalera FentMates de les fitxes: un color per a cada curs.
% Els colors són els de la pàgina inicial (LEVELS a assets/js/materials.js);
% el logo ve de _eines/logo/pdf/ (python3 _eines/logo/genera.py).
%
% Ús, al final del preàmbul de la unitat:
%   \input{../comu/fentmates.tex}
%   \fmcurs{1eso}
% i, al començament del document:
%   \capcalera{Matemàtiques}{1r ESO}{Fitxa d'exercicis · Unitat 2}{Divisibilitat}
% Opcional abans de \input: \def\fmneed{14} (espai mínim abans de cada secció).
% =====================================================================
\makeatletter
\def\fm@logodir{../../logo/pdf}
\providecommand{\fmneed}{10}
\def\fm@cursnom{}

% clau del curs, color i color del text damunt del color
\newcommand{\fm@curs}[3]{\@namedef{fm@color@#1}{#2}\@namedef{fm@text@#1}{#3}}
\fm@curs{1eso}{E6B84A}{navy}
\fm@curs{2eso}{3B82F6}{white}
\fm@curs{3eso}{FF7A50}{white}
\fm@curs{4eso-a}{A78BFA}{white}
\fm@curs{4eso-b}{6D28D9}{white}
\fm@curs{1batx-mi}{9B1C1C}{white}
\fm@curs{1batx-ccss}{EF6B6B}{white}
\fm@curs{2batx-mii}{1E6B47}{white}
\fm@curs{2batx-ccss}{2BB3A3}{white}

\newcommand{\fmcurs}[1]{%
  \@ifundefined{fm@color@#1}{\PackageError{fentmates}{Curs desconegut: #1}{}}{}%
  \edef\fm@tmp{\noexpand\definecolor{curs}{HTML}{\@nameuse{fm@color@#1}}%
               \noexpand\colorlet{curstext}{\@nameuse{fm@text@#1}}}\fm@tmp
  \colorlet{cursfosc}{curs!72!black}%
  \def\fm@logo{\fm@logodir/vertical-#1.pdf}%
  \def\fm@marca{\fm@logodir/marca-#1.pdf}}

% ---------- capçalera de la primera pàgina ----------
% #1 matèria (petit), #2 curs, #3 tipus de document i unitat, #4 títol
\newcommand{\capcalera}[4]{%
  \gdef\fm@cursnom{#2}%
  \par\noindent
  \begin{tcolorbox}[colback=curs,colframe=curs,boxrule=0pt,arc=2.2mm,width=3.5cm,height=2.15cm,
      valign=center,left=3mm,right=2mm,top=1mm,bottom=1mm,nobeforeafter,box align=center]
    \color{curstext}{\footnotesize #1}\par\vspace{1pt}{\fontsize{19}{21}\selectfont\bfseries #2\par}%
  \end{tcolorbox}\hspace{5mm}%
  \begin{minipage}[c]{\dimexpr\linewidth-3.5cm-5mm-2.9cm\relax}
    {\color{cursfosc}\small\bfseries\MakeUppercase{#3}}\par\vspace{2pt}
    {\color{navy}\fontsize{20}{23}\selectfont\bfseries #4\par}
  \end{minipage}\hfill
  \begin{minipage}[c]{2.6cm}\raggedleft\includegraphics[width=2.6cm]{\fm@logo}\end{minipage}\par
  \vspace{2pt}}

% ---------- barres de secció del color del curs ----------
\renewcommand{\seccio}[2]{\par\needspace{\fmneed\baselineskip}\bigskip
  \begin{tcolorbox}[colback=curs,colframe=curs,coltext=curstext,arc=2mm,boxrule=0pt,left=3mm,right=3mm,top=1.2mm,bottom=1.2mm]
  {\large\bfseries #1\quad #2}
  \end{tcolorbox}\vspace{-2pt}}

% ---------- estrelles de dificultat del color del curs (a Batxillerat no se'n posen) ----------
\@ifundefined{estrelles}{}{%
  \renewcommand{\estrelles}[1]{{\color{curs}\ifcase#1\or\faStar\faStar[regular]\faStar[regular]\or\faStar\faStar\faStar[regular]\or\faStar\faStar\faStar\fi}}}

% ---------- capçalera i peu de les pàgines ----------
% primera pàgina: logos institucionals al peu
\fancypagestyle{primera}{\fancyhf{}\renewcommand{\headrulewidth}{0pt}
  \fancyfoot[L]{\raisebox{-0.32\height}{\includegraphics[height=0.78cm]{logo-gva.png}}\hspace{5mm}%
    \raisebox{-0.32\height}{\includegraphics[height=0.78cm]{logo-vermellar.png}}}
  \fancyfoot[R]{\scriptsize\color{muted}\href{https://creativecommons.org/licenses/by-nc-sa/4.0/deed.ca}{\ccbyncsa}\ \ CC BY-NC-SA 4.0 · fentmates.org\qquad\small\thepage}}
% la resta: el símbol a l'esquerra i el curs a la dreta, amb el títol al centre com abans
\fancyhead[L]{\raisebox{-0.3\height}{\includegraphics[height=0.62cm]{\fm@marca}}}
\fancyhead[R]{\small\bfseries\color{cursfosc}\fm@cursnom}
\renewcommand{\headrule}{\color{curs!45}\hrule height \headrulewidth}
\makeatother

%   \fmcurs{1eso}
% i, al començament del document:
%   \capcalera{Matemàtiques}{1r ESO}{Fitxa d'exercicis · Unitat 2}{Divisibilitat}
% Opcional abans de \input: \def\fmneed{14} (espai mínim abans de cada secció).
% =====================================================================
\makeatletter
\def\fm@logodir{../../logo/pdf}
\providecommand{\fmneed}{10}
\def\fm@cursnom{}

% clau del curs, color i color del text damunt del color
\newcommand{\fm@curs}[3]{\@namedef{fm@color@#1}{#2}\@namedef{fm@text@#1}{#3}}
\fm@curs{1eso}{E6B84A}{navy}
\fm@curs{2eso}{3B82F6}{white}
\fm@curs{3eso}{FF7A50}{white}
\fm@curs{4eso-a}{A78BFA}{white}
\fm@curs{4eso-b}{6D28D9}{white}
\fm@curs{1batx-mi}{9B1C1C}{white}
\fm@curs{1batx-ccss}{EF6B6B}{white}
\fm@curs{2batx-mii}{1E6B47}{white}
\fm@curs{2batx-ccss}{2BB3A3}{white}

\newcommand{\fmcurs}[1]{%
  \@ifundefined{fm@color@#1}{\PackageError{fentmates}{Curs desconegut: #1}{}}{}%
  \edef\fm@tmp{\noexpand\definecolor{curs}{HTML}{\@nameuse{fm@color@#1}}%
               \noexpand\colorlet{curstext}{\@nameuse{fm@text@#1}}}\fm@tmp
  \colorlet{cursfosc}{curs!72!black}%
  \def\fm@logo{\fm@logodir/vertical-#1.pdf}%
  \def\fm@marca{\fm@logodir/marca-#1.pdf}}

% ---------- capçalera de la primera pàgina ----------
% #1 matèria (petit), #2 curs, #3 tipus de document i unitat, #4 títol
\newcommand{\capcalera}[4]{%
  \gdef\fm@cursnom{#2}%
  \par\noindent
  \begin{tcolorbox}[colback=curs,colframe=curs,boxrule=0pt,arc=2.2mm,width=3.5cm,height=2.15cm,
      valign=center,left=3mm,right=2mm,top=1mm,bottom=1mm,nobeforeafter,box align=center]
    \color{curstext}{\footnotesize #1}\par\vspace{1pt}{\fontsize{19}{21}\selectfont\bfseries #2\par}%
  \end{tcolorbox}\hspace{5mm}%
  \begin{minipage}[c]{\dimexpr\linewidth-3.5cm-5mm-2.9cm\relax}
    {\color{cursfosc}\small\bfseries\MakeUppercase{#3}}\par\vspace{2pt}
    {\color{navy}\fontsize{20}{23}\selectfont\bfseries #4\par}
  \end{minipage}\hfill
  \begin{minipage}[c]{2.6cm}\raggedleft\includegraphics[width=2.6cm]{\fm@logo}\end{minipage}\par
  \vspace{2pt}}

% ---------- barres de secció del color del curs ----------
\renewcommand{\seccio}[2]{\par\needspace{\fmneed\baselineskip}\bigskip
  \begin{tcolorbox}[colback=curs,colframe=curs,coltext=curstext,arc=2mm,boxrule=0pt,left=3mm,right=3mm,top=1.2mm,bottom=1.2mm]
  {\large\bfseries #1\quad #2}
  \end{tcolorbox}\vspace{-2pt}}

% ---------- estrelles de dificultat del color del curs (a Batxillerat no se'n posen) ----------
\@ifundefined{estrelles}{}{%
  \renewcommand{\estrelles}[1]{{\color{curs}\ifcase#1\or\faStar\faStar[regular]\faStar[regular]\or\faStar\faStar\faStar[regular]\or\faStar\faStar\faStar\fi}}}

% ---------- capçalera i peu de les pàgines ----------
% primera pàgina: logos institucionals al peu
\fancypagestyle{primera}{\fancyhf{}\renewcommand{\headrulewidth}{0pt}
  \fancyfoot[L]{\raisebox{-0.32\height}{\includegraphics[height=0.78cm]{logo-gva.png}}\hspace{5mm}%
    \raisebox{-0.32\height}{\includegraphics[height=0.78cm]{logo-vermellar.png}}}
  \fancyfoot[R]{\scriptsize\color{muted}\href{https://creativecommons.org/licenses/by-nc-sa/4.0/deed.ca}{\ccbyncsa}\ \ CC BY-NC-SA 4.0 · fentmates.org\qquad\small\thepage}}
% la resta: el símbol a l'esquerra i el curs a la dreta, amb el títol al centre com abans
\fancyhead[L]{\raisebox{-0.3\height}{\includegraphics[height=0.62cm]{\fm@marca}}}
\fancyhead[R]{\small\bfseries\color{cursfosc}\fm@cursnom}
\renewcommand{\headrule}{\color{curs!45}\hrule height \headrulewidth}
\makeatother

%   \fmcurs{1eso}
% i, al començament del document:
%   \capcalera{Matemàtiques}{1r ESO}{Fitxa d'exercicis · Unitat 2}{Divisibilitat}
% Opcional abans de \input: \def\fmneed{14} (espai mínim abans de cada secció).
% =====================================================================
\makeatletter
\def\fm@logodir{../../logo/pdf}
\providecommand{\fmneed}{10}
\def\fm@cursnom{}

% clau del curs, color i color del text damunt del color
\newcommand{\fm@curs}[3]{\@namedef{fm@color@#1}{#2}\@namedef{fm@text@#1}{#3}}
\fm@curs{1eso}{E6B84A}{navy}
\fm@curs{2eso}{3B82F6}{white}
\fm@curs{3eso}{FF7A50}{white}
\fm@curs{4eso-a}{A78BFA}{white}
\fm@curs{4eso-b}{6D28D9}{white}
\fm@curs{1batx-mi}{9B1C1C}{white}
\fm@curs{1batx-ccss}{EF6B6B}{white}
\fm@curs{2batx-mii}{1E6B47}{white}
\fm@curs{2batx-ccss}{2BB3A3}{white}

\newcommand{\fmcurs}[1]{%
  \@ifundefined{fm@color@#1}{\PackageError{fentmates}{Curs desconegut: #1}{}}{}%
  \edef\fm@tmp{\noexpand\definecolor{curs}{HTML}{\@nameuse{fm@color@#1}}%
               \noexpand\colorlet{curstext}{\@nameuse{fm@text@#1}}}\fm@tmp
  \colorlet{cursfosc}{curs!72!black}%
  \def\fm@logo{\fm@logodir/vertical-#1.pdf}%
  \def\fm@marca{\fm@logodir/marca-#1.pdf}}

% ---------- capçalera de la primera pàgina ----------
% #1 matèria (petit), #2 curs, #3 tipus de document i unitat, #4 títol
\newcommand{\capcalera}[4]{%
  \gdef\fm@cursnom{#2}%
  \par\noindent
  \begin{tcolorbox}[colback=curs,colframe=curs,boxrule=0pt,arc=2.2mm,width=3.5cm,height=2.15cm,
      valign=center,left=3mm,right=2mm,top=1mm,bottom=1mm,nobeforeafter,box align=center]
    \color{curstext}{\footnotesize #1}\par\vspace{1pt}{\fontsize{19}{21}\selectfont\bfseries #2\par}%
  \end{tcolorbox}\hspace{5mm}%
  \begin{minipage}[c]{\dimexpr\linewidth-3.5cm-5mm-2.9cm\relax}
    {\color{cursfosc}\small\bfseries\MakeUppercase{#3}}\par\vspace{2pt}
    {\color{navy}\fontsize{20}{23}\selectfont\bfseries #4\par}
  \end{minipage}\hfill
  \begin{minipage}[c]{2.6cm}\raggedleft\includegraphics[width=2.6cm]{\fm@logo}\end{minipage}\par
  \vspace{2pt}}

% ---------- barres de secció del color del curs ----------
\renewcommand{\seccio}[2]{\par\needspace{\fmneed\baselineskip}\bigskip
  \begin{tcolorbox}[colback=curs,colframe=curs,coltext=curstext,arc=2mm,boxrule=0pt,left=3mm,right=3mm,top=1.2mm,bottom=1.2mm]
  {\large\bfseries #1\quad #2}
  \end{tcolorbox}\vspace{-2pt}}

% ---------- estrelles de dificultat del color del curs (a Batxillerat no se'n posen) ----------
\@ifundefined{estrelles}{}{%
  \renewcommand{\estrelles}[1]{{\color{curs}\ifcase#1\or\faStar\faStar[regular]\faStar[regular]\or\faStar\faStar\faStar[regular]\or\faStar\faStar\faStar\fi}}}

% ---------- capçalera i peu de les pàgines ----------
% primera pàgina: logos institucionals al peu
\fancypagestyle{primera}{\fancyhf{}\renewcommand{\headrulewidth}{0pt}
  \fancyfoot[L]{\raisebox{-0.32\height}{\includegraphics[height=0.78cm]{logo-gva.png}}\hspace{5mm}%
    \raisebox{-0.32\height}{\includegraphics[height=0.78cm]{logo-vermellar.png}}}
  \fancyfoot[R]{\scriptsize\color{muted}\href{https://creativecommons.org/licenses/by-nc-sa/4.0/deed.ca}{\ccbyncsa}\ \ CC BY-NC-SA 4.0 · fentmates.org\qquad\small\thepage}}
% la resta: el símbol a l'esquerra i el curs a la dreta, amb el títol al centre com abans
\fancyhead[L]{\raisebox{-0.3\height}{\includegraphics[height=0.62cm]{\fm@marca}}}
\fancyhead[R]{\small\bfseries\color{cursfosc}\fm@cursnom}
\renewcommand{\headrule}{\color{curs!45}\hrule height \headrulewidth}
\makeatother

\fmcurs{4eso-a}

\begin{document}
\thispagestyle{primera}

\capcalera{Matemàtiques A}{4t ESO}{Fitxa d'exercicis · Unitat 1}{\raggedright Nombres naturals, enters i fraccionaris}
\vspace{10pt}
Nom i cognoms: \rule{8.4cm}{0.4pt}\hfill Grup: \rule{1.6cm}{0.4pt}\hfill Data: \rule{2.2cm}{0.4pt}
\vspace{8pt}

\begin{tcolorbox}[colback=soft,colframe=soft,arc=2mm,left=3mm,right=3mm,top=1.5mm,bottom=1.5mm]
\small Exercicis graduats per nivells de dificultat:\quad \estrelles{1}\ bàsic \qquad \estrelles{2}\ mitjà \qquad \estrelles{3}\ avançat.\\[2pt] Resol-los al quadern mostrant tots els passos. Fes els càlculs sense calculadora.
\end{tcolorbox}

% =====================================================
\seccio{1}{Nombres naturals. Operacions combinades}

\ex{1} Calcula:\begin{tasks}(6)
  \task $3^{5}$
  \task $10^{6}$
  \task $2^{3}\cdot 5^{2}$
  \task $\sqrt{144}$
  \task $\sqrt{225}$
  \task $17^{0}$
\end{tasks}

\ex{1} Calcula el quocient i el residu de cada divisió i fes la prova ($D=d\cdot q+r$):\begin{tasks}(2)
  \task $1250:37$
  \task $5000:123$
\end{tasks}

\ex{2} Calcula pas a pas:\begin{tasks}[after-item-skip=6pt](2)
  \task $25-3\cdot \left[12-2\cdot \left(7-4\right)\right]$
  \task $6+48:\left[20-4\cdot \left(9-6\right)\right]$
  \task $60:\left[3\cdot \left(9-5\right)+2\cdot \left(7-3\right)\right]$
  \task $2\cdot \left(40-3\cdot \left[5+2\cdot \left(9-6\right)\right]\right)+7$
  \task $5\cdot \left[18-\left(4+2\cdot 3\right)\right]-4\cdot \left[10-\left(9-2\cdot 3\right)\right]$
  \task $2^{3}\cdot \left[15-\left(\sqrt{49}+2\right)\right]+36:3^{2}$
\end{tasks}\recorda{Primer els parèntesis i claudàtors (de dins cap a fora); després les potències i les arrels; després les multiplicacions i divisions (d'esquerra a dreta) i, al final, les sumes i restes.}

\ex{2} Comprem 3 samarretes a 18\,€ i 2 pantalons a 35\,€, i ens fan un descompte de 15\,€ sobre el total. Paguem amb dos bitllets de 100\,€. Escriu l'operació en una sola expressió i calcula el canvi.

\ex{3} Un alumne ha fet $36:6\cdot2+4^2=36:12+16=3+16=19$. Explica l'error i fes-ho bé.

\ex{3} Un hort quadrat té una superfície de 1225\,m². El volem tancar amb 3 voltes de filferro, deixant una porta de 5\,m. El filferro es ven en rotlles de 50\,m a 18\,€ el rotlle. Quants rotlles cal comprar? Quant costaran?


% =====================================================
\seccio{2}{Nombres enters. Valor absolut}

\ex{1} Ordena de menor a major: $\ {-15},\ 8,\ {-3},\ 0,\ {-21},\ 5,\ {-8},\ 12$.

\ex{1} Calcula:\begin{tasks}[after-item-skip=4pt](3)
  \task $\left\lvert -15\right\rvert +\left\lvert 7\right\rvert $
  \task $\left\lvert -4-9\right\rvert $
  \task $\left\lvert -4\right\rvert -\left\lvert 9\right\rvert $
  \task $-\left\lvert -12\right\rvert $
  \task $\left\lvert 3-11\right\rvert \cdot \left\lvert -2\right\rvert $
  \task $\left\lvert -20\right\rvert :\left\lvert -5\right\rvert -\left\lvert -3\right\rvert $
\end{tasks}\recorda{El valor absolut $\lvert a\rvert $ és la distància de $a$ al 0: sempre és positiu o zero. Si dins hi ha operacions, primer les fem i després traiem el valor absolut.}

\ex{2} Escriu tots els nombres enters $x$ que compleixen:\begin{tasks}(4)
  \task $\lvert x\rvert =9$
  \task $\lvert x\rvert \le3$
  \task $-4\le x<2$
  \task $\lvert x-2\rvert =5$
\end{tasks}

\ex{3} Calcula pas a pas:\begin{tasks}[after-item-skip=6pt](2)
  \task $\left\lvert 4-9\right\rvert -\left\lvert 2-3\cdot 5\right\rvert +\left\lvert -6\right\rvert \cdot \left(-2\right)$
  \task $\left(3-8\right)\cdot 2+\left\lvert 6-4\cdot 3\right\rvert \cdot \left\lvert -1-2\right\rvert $
  \task $7-\left\lvert -3-\left\lvert 2-9\right\rvert \right\rvert $
  \task $\left\lvert \left(-2\right)\cdot \left(-5\right)-\left\lvert 3-15\right\rvert \right\rvert -\left\lvert -4\right\rvert \cdot \left(1-3\right)$
\end{tasks}

\ex{2} A les 7 del matí, en una ciutat de Sibèria, el termòmetre marcava $-27$\,°C. A migdia la temperatura havia pujat 19 graus i a mitjanit havia baixat 14 graus respecte de migdia. a) Quina temperatura feia a migdia i a mitjanit? b) Quina diferència hi ha entre la temperatura més alta i la més baixa d'aquestes tres?

\ex{2} Arquimedes va nàixer l'any 287 aC i va morir l'any 212 aC. Quants anys va viure?


% =====================================================
\seccio{3}{Operacions combinades amb enters}

\ex{1} Calcula:\begin{tasks}(3)
  \task $-12+7-\left(-5\right)$
  \task $\left(-8\right)\cdot \left(-6\right):\left(-4\right)$
  \task $\left(-72\right):\left(-9\right)\cdot \left(-2\right)$
  \task $-15-\left(-3\right)\cdot 4$
  \task $\left(-2\right)\cdot \left(-3\right)\cdot \left(-5\right)$
  \task $-36:\left(-4\right)-9$
\end{tasks}\recorda{Regla dels signes: $(+)\cdot(+)=+$, \ $(-)\cdot(-)=+$, \ $(+)\cdot(-)=-$, \ $(-)\cdot(+)=-$. També per a la divisió.}

\ex{2} Calcula pas a pas:\begin{tasks}[after-item-skip=6pt](2)
  \task $12-4\cdot \left[6-3\cdot \left(5-8\right)\right]$
  \task $-3+24:\left[-2-2\cdot \left(1-6\right)\right]$
  \task $5\cdot \left(-4+9\right)-\left[7-\left(-3\right)\cdot \left(2-6\right)\right]$
  \task $\left(-2\right)\cdot \left[8-\left(3-5\right)\cdot \left(-4\right)\right]-\left(-6\right):\left(-2\right)$
  \task $40:\left[-12+4\cdot \left(7-9\right)\cdot \left(-1\right)\right]$
  \task $-\left[6-\left(2-11\right)\right]+3\cdot \left[-4+\left(5-1\right)\cdot 2\right]$
\end{tasks}

\ex{2} Calcula pas a pas (amb potències i arrels):\begin{tasks}[after-item-skip=6pt](2)
  \task $\left(-3\right)^{3}+\left(-2\right)^{4}-\left(-1\right)^{9}$
  \task $-4^{2}+\left(-4\right)^{2}-\sqrt{36}\cdot \left(-2\right)^{3}$
  \task $\left[\left(-5\right)^{2}-\sqrt{100-36}\cdot 3\right]:\left(-1\right)^{7}$
  \task $2\cdot \left(-3\right)^{2}-\left[\left(-2\right)^{3}+\sqrt{81}\right]^{2}$
  \task $\left(-2\right)^{5}:\left(-2\right)^{3}-3\cdot \left(-1\right)^{20}+\sqrt{4^{2}+3^{2}}$
\end{tasks}\recorda{$(-4)^2=16$, però $-4^2=-16$: sense parèntesi, el signe no forma part de la base.}

\ex{2} Un compte bancari té un saldo de 420\,€. Es paguen 3 rebuts de 85\,€ i 2 de 120\,€ i, després, s'ingressen dues nòmines de 150\,€. Escriu-ho en una sola expressió i calcula el saldo final. Ha arribat a estar en negatiu?

\ex{3} Un ascensor és a la planta 3. Baixa 5 plantes, després puja el doble de les que ha baixat i, finalment, baixa 9 plantes. A quina planta està? Si entre planta i planta hi ha 3\,m, quants metres ha recorregut en total?

\ex{3} Les temperatures mínimes d'una setmana han sigut $-6$, $3$, $-2$, $-8$, $1$, $-4$ i $2$\,°C. Calcula la temperatura mitjana i la diferència entre la màxima i la mínima.

\ex{3} Troba dos nombres enters que multiplicats donen $-24$ i sumats donen $-2$.


% =====================================================
\seccio{4}{Fraccions. Fraccions equivalents}

\ex{1} Simplifica fins a obtindre la fracció irreductible:\begin{tasks}[after-item-skip=8pt](6)
  \task $\frac{84}{140}$
  \task $\frac{126}{315}$
  \task $\frac{210}{462}$
  \task $\frac{360}{840}$
  \task $\frac{156}{234}$
  \task $\frac{275}{1100}$
\end{tasks}\recorda{Per a arribar directament a la fracció irreductible, dividim el numerador i el denominador entre el seu MCD.}

\ex{1} Calcula:\begin{tasks}(1)
  \task Els $\frac{2}{3}$ dels $\frac{3}{5}$ de 450\,€.
  \task Els $\frac{3}{4}$ de la meitat de 96 alumnes.
  \task Els $\frac{5}{6}$ d\'un terç de 720 litres.
\end{tasks}

\ex{2} Comprova si són equivalents:\begin{tasks}[after-item-skip=6pt](3)
  \task $\frac{21}{35}$ i $\frac{33}{55}$
  \task $\frac{18}{24}$ i $\frac{27}{32}$
  \task $\frac{-14}{49}$ i $\frac{6}{-21}$
\end{tasks}

\ex{2} Troba el valor de $x$ perquè les fraccions siguen equivalents:\begin{tasks}[after-item-skip=6pt](3)
  \task $\frac{x}{18}=\frac{35}{42}$
  \task $\frac{12}{x}=-\frac{4}{7}$
  \task $\frac{-15}{25}=\frac{x}{-40}$
\end{tasks}

\ex{2} Quina fracció del dia són 16 hores? I 45 minuts? I 2 h 40 min? Dona-les irreductibles.

\ex{2} Ordena de menor a major i situa-les aproximadament en la recta: $\ \frac{5}{6},\ -\frac{3}{4},\ \frac{7}{9},\ -\frac{2}{3},\ \frac{13}{18},\ 0$.

\ex{3} Escriu una fracció compresa entre $\frac{3}{7}$ i $\frac{4}{7}$, i una altra entre $\frac{2}{5}$ i $\frac{1}{2}$. Explica com ho has fet.


% =====================================================
\seccio{5}{Operacions amb fraccions}

\ex{1} Calcula i simplifica el resultat:\begin{tasks}[after-item-skip=8pt](4)
  \task $\frac{5}{12}+\frac{7}{18}$
  \task $\frac{11}{15}-\frac{7}{10}$
  \task $3-\frac{11}{4}+\frac{5}{6}$
  \task $-\frac{7}{20}+\frac{5}{8}$
\end{tasks}\recorda{Per a sumar o restar, primer reduïm a comú denominador (amb el mcm dels denominadors). Mai no sumem els denominadors!}

\ex{1} Calcula i simplifica el resultat:\begin{tasks}[after-item-skip=8pt](4)
  \task $\frac{14}{15}\cdot \frac{25}{21}$
  \task $-\frac{9}{16}\cdot \frac{8}{27}$
  \task $\frac{12}{35}:\frac{18}{25}$
  \task $\left(-\frac{5}{6}\right):\left(-\frac{10}{9}\right)$
\end{tasks}\recorda{$\frac{a}{b}\cdot\frac{c}{d}=\frac{a\cdot c}{b\cdot d}$ \qquad $\frac{a}{b}:\frac{c}{d}=\frac{a\cdot d}{b\cdot c}$. Simplifica abans de multiplicar i els números seran més menuts.}

\ex{2} Troba la fracció resultant:\begin{tasks}[after-item-skip=6pt](4)
  \task $\frac{2}{3}$ de $\frac{9}{10}$
  \task $la meitat$ de $\frac{4}{7}$
  \task $\frac{3}{5}$ de $\frac{5}{12}$
  \task $\frac{7}{8}$ de $\frac{4}{21}$
\end{tasks}

\ex{2} Calcula (compte amb els signes):\begin{tasks}[after-item-skip=8pt](2)
  \task $-\frac{5}{6}-\left(-\frac{3}{4}\right)+-\frac{1}{3}$
  \task $\left(-\frac{2}{3}\right)\cdot \frac{9}{4}\cdot \left(-\frac{8}{15}\right)$
  \task $\left(-\frac{4}{9}\right):\left(-\frac{2}{3}\right)\cdot \left(-6\right)$
  \task $-2-\left(-\frac{5}{4}\right)+\frac{7}{6}$
\end{tasks}

\ex{2} Escriu l'oposat i l'invers de cada nombre:\begin{tasks}(4)
  \task $\frac{3}{8}$
  \task $-\frac{5}{2}$
  \task $-6$
  \task $1+\frac{1}{3}$
\end{tasks}


% =====================================================
\seccio{6}{Operacions combinades amb fraccions}

\ex{1} Calcula i compara els resultats de a) i b), i de c) i d):\begin{tasks}[after-item-skip=8pt](4)
  \task $\frac{2}{3}+\frac{1}{2}\cdot \frac{4}{9}$
  \task $\left(\frac{2}{3}+\frac{1}{2}\right)\cdot \frac{4}{9}$
  \task $\frac{3}{4}-\frac{1}{4}:\frac{3}{8}$
  \task $\left(\frac{3}{4}-\frac{1}{4}\right):\frac{3}{8}$
\end{tasks}

\ex{2} Calcula pas a pas:\begin{tasks}[after-item-skip=10pt](2)
  \task $\frac{3}{4}-\left[\frac{1}{2}-\left(\frac{2}{5}-\frac{1}{10}\right)\right]$
  \task $2-\left[\frac{1}{3}+\left(\frac{5}{6}-1\right)\right]+\frac{1}{4}$
  \task $1-\left(\frac{2}{3}+\frac{1}{4}-\frac{1}{2}\right):\frac{5}{6}$
  \task $\frac{4}{5}\cdot \left[\frac{1}{2}-\left(\frac{3}{4}-\frac{1}{8}\right)\right]$
  \task $\left[\frac{2}{3}\cdot \left(\frac{1}{4}+\frac{1}{2}\right)\right]:\left[\left(-3\right)\cdot \left(\frac{1}{6}-\frac{1}{2}\right)\right]$
  \task $\left(\frac{5}{6}-\frac{2}{3}-\frac{1}{12}\right):\left(1-\frac{3}{4}\right)$
\end{tasks}

\ex{2} Calcula pas a pas (amb potències):\begin{tasks}[after-item-skip=8pt](3)
  \task $\left(\frac{3}{2}\right)^{2}-\frac{5}{4}\cdot \left(1-\frac{1}{5}\right)$
  \task $\left(-\frac{2}{3}\right)^{3}:\frac{4}{9}+\frac{1}{2}$
  \task $\left[1-\left(\frac{1}{2}\right)^{2}\right]^{2}\cdot \frac{16}{3}$
\end{tasks}\recorda{Per a elevar una fracció, elevem el numerador i el denominador: $\left(-\frac{2}{3}\right)^3=-\frac{2^3}{3^3}=-\frac{8}{27}$.}

\ex{3} Calcula (primer el numerador i el denominador, i després divideix):\begin{tasks}[after-item-skip=10pt](4)
  \task $\dfrac{\frac{2}{3}-\frac{1}{4}}{\frac{5}{6}+\frac{1}{3}}$
  \task $\dfrac{1-\frac{3}{5}\cdot \frac{5}{6}}{\left(\frac{1}{2}\right)^{2}+\frac{1}{3}}$
  \task $\dfrac{\left(-2\right)\cdot \left(\frac{1}{4}-\frac{2}{3}\right)}{3\cdot \left(\frac{5}{6}-\frac{1}{2}\right)}$
  \task $\dfrac{\frac{3}{4}:\frac{9}{8}-\frac{1}{3}}{1-\frac{4}{9}}$
\end{tasks}

\ex{3} Calcula pas a pas:\begin{tasks}[after-item-skip=8pt](1)
  \task $\left[\frac{3}{4}-\left(\frac{1}{2}\right)^{2}\right]^{2}:\left(-\frac{1}{2}\right)^{3}+\frac{5}{3}\cdot \frac{6}{5}$
  \task $\frac{1}{2}-\frac{2}{3}\cdot \left[\frac{3}{4}-\left(\frac{1}{2}-2\right)^{2}:3\right]$
\end{tasks}


% =====================================================
\seccio{7}{Problemes amb fraccions}

\ex{1} Els $\frac{4}{9}$ dels alumnes d'un institut, que són 352, venen amb autobús. Quants alumnes té l'institut?

\ex{2} Un mòbil costa 480\,€. Paguem $\frac{1}{4}$ d'entrada i la resta en 9 quotes iguals. Quant paguem en cada quota? Quina fracció del preu és cada quota?

\ex{2} Es reparteixen 1800\,€ entre tres germans: el primer rep $\frac{2}{5}$ del total, el segon $\frac{3}{4}$ del que queda i el tercer la resta. Quant rep cada un? Quina fracció del total rep el tercer?

\ex{2} En un camp de 2400\,m² es planten tomaques en $\frac{3}{8}$ de la superfície i creïlles en $\frac{2}{5}$ de la resta. La resta és per a pebrots. Quants m² ocupa cada cultiu? Quina fracció del camp ocupen els pebrots?

\ex{3} Una aixeta omple $\frac{1}{6}$ d'un dipòsit cada hora i una altra, $\frac{1}{4}$. Si les obrim totes dues alhora, quina fracció del dipòsit s'omple en una hora? Quant de temps tarden a omplir-lo?

\ex{3} L'Anna gasta $\frac{2}{7}$ del sou en el lloguer i $\frac{1}{3}$ del que li queda en menjar. Encara li queden 600\,€. Quant cobra?

\ex{3} En un viatge, el primer dia fem $\frac{2}{5}$ del recorregut i el segon dia, $\frac{3}{4}$ del que falta. Encara queden 90\,km. Quants quilòmetres té el viatge?

\ex{3} Un pot ple de mel pesa 950\,g. Quan n'hem gastat $\frac{3}{5}$ de la mel, pesa 470\,g. Quant pesa el pot buit?

\ex{3} D'un dipòsit ple, el dilluns se'n gasta $\frac{1}{3}$; el dimarts, $\frac{1}{4}$ del que quedava, i el dimecres, $\frac{2}{5}$ del que quedava. Al final queden 360 litres. Quina és la capacitat del dipòsit?


% =====================================================
\seccio{8}{Potències amb base entera i racional}

\ex{1} Calcula:\begin{tasks}(6)
  \task $\left(-4\right)^{3}$
  \task $-4^{2}$
  \task $\left(-1\right)^{23}$
  \task $-\left(-2\right)^{4}$
  \task $\left(-10\right)^{5}$
  \task $\left(-9\right)^{0}$
\end{tasks}

\ex{1} Calcula:\begin{tasks}[after-item-skip=8pt](6)
  \task $\left(\frac{3}{4}\right)^{3}$
  \task $\left(-\frac{2}{3}\right)^{4}$
  \task $\left(-\frac{1}{5}\right)^{3}$
  \task $-\left(\frac{3}{2}\right)^{2}$
  \task $\left(-\frac{7}{8}\right)^{0}$
  \task $\left(-\frac{1}{10}\right)^{5}$
\end{tasks}\recorda{$\left(\frac{a}{b}\right)^n=\frac{a^n}{b^n}$. Si la base és negativa, el resultat és positiu amb exponent parell i negatiu amb exponent senar.}

\ex{2} Sense fer el càlcul, digues quin signe té el resultat:\begin{tasks}(5)
  \task $(-6)^{9}$
  \task $(-11)^{12}$
  \task $-3^{8}$
  \task $\left(-\frac{2}{5}\right)^{11}$
  \task $-(-1)^{100}$
\end{tasks}

\ex{2} Calcula:\begin{tasks}[after-item-skip=8pt](2)
  \task $\left(-3\right)^{3}-\left(-3\right)^{2}-3^{2}$
  \task $\left(-1\right)^{15}-\left(-1\right)^{16}+\left(-2\right)^{3}$
  \task $\left(\frac{2}{3}\right)^{2}-\left(\frac{1}{3}\right)^{2}$
  \task $\left(-\frac{1}{2}\right)^{3}\cdot 8-\left(\frac{3}{4}\right)^{2}:\frac{9}{8}$
\end{tasks}

\ex{3} Deixem caure una pilota des de 16\,m d'altura. Cada vegada que bota, puja $\frac{3}{4}$ de l'altura anterior. a) A quina altura arriba després del quart bot? b) Quants metres ha recorregut en total quan toca terra per tercera vegada?


% =====================================================
\seccio{9}{Propietats de les potències}

\ex{1} Expressa com una sola potència:\begin{tasks}(3)
  \task $5^{3}\cdot 5^{4}$
  \task $2^{9}:2^{4}$
  \task $\left(3^{2}\right)^{5}$
  \task $4^{3}\cdot 5^{3}$
  \task $18^{6}:6^{6}$
  \task $\left(-7\right)^{4}\cdot \left(-7\right)^{3}$
\end{tasks}\recorda{Mateixa base: $a^m\cdot a^n=a^{m+n}$, \ $a^m:a^n=a^{m-n}$. \ Potència de potència: $(a^m)^n=a^{m\cdot n}$. \ Mateix exponent: $a^n\cdot b^n=(a\cdot b)^n$, \ $a^n:b^n=(a:b)^n$.}

\ex{1} Expressa com una sola potència:\begin{tasks}[after-item-skip=8pt](2)
  \task $\left(\frac{2}{5}\right)^3\cdot\left(\frac{2}{5}\right)^4$
  \task $\left(\frac{3}{4}\right)^8:\left(\frac{3}{4}\right)^5$
  \task $\left[\left(-\frac{1}{2}\right)^3\right]^2$
  \task $\left(\frac{2}{3}\right)^4\cdot\left(\frac{9}{4}\right)^4$
\end{tasks}

\ex{2} Simplifica aplicant les propietats i calcula:\begin{tasks}[after-item-skip=8pt](3)
  \task $\left(-2\right)^{7}:\left(-2\right)^{4}$
  \task $\left[\left(-1\right)^{3}\right]^{5}\cdot \left(-1\right)^{2}$
  \task $2^{5}\cdot 2^{3}:2^{6}$
  \task $\left(\frac{2}{3}\right)^{5}\cdot \left(\frac{2}{3}\right)^{2}:\left(\frac{2}{3}\right)^{4}$
  \task $\left(-\frac{1}{3}\right)^{6}:\left(-\frac{1}{3}\right)^{4}$
\end{tasks}

\ex{2} Vertader o fals? Corregeix les falses:\begin{tasks}(3)
  \task $2^3+2^4=2^7$
  \task $3^2\cdot3^3=3^5$
  \task $(2+3)^2=2^2+3^2$
  \task $(3^2)^3=3^5$
  \task $5^0=0$
  \task $6^4:3^4=2^4$
\end{tasks}

\ex{3} Calcula descomponent les bases en factors primers:\begin{tasks}[after-item-skip=8pt](3)
  \task $\dfrac{10^4\cdot2^3}{5^3\cdot4^5}$
  \task $9^3\cdot27:3^7$
  \task $\dfrac{12^3\cdot6^2}{18^2\cdot2^5}$
\end{tasks}

\ex{3} Calcula de la manera més ràpida: $\ \left(\frac{4}{9}\right)^2\cdot\left(\frac{3}{2}\right)^4$


% =====================================================
\seccio{10}{Potències d'exponent negatiu}

\ex{1} Calcula el valor d'aquestes potències:\begin{tasks}[after-item-skip=8pt](4)
  \task $5^{-2}$
  \task $2^{-5}$
  \task $10^{-3}$
  \task $\left(-3\right)^{-2}$
  \task $\left(-2\right)^{-3}$
  \task $\left(\frac{1}{4}\right)^{-2}$
  \task $\left(\frac{2}{5}\right)^{-3}$
  \task $\left(-\frac{3}{2}\right)^{-1}$
\end{tasks}\recorda{$a^0=1$, \ $a^{-n}=\dfrac{1}{a^n}$ \ i \ $\left(\dfrac{a}{b}\right)^{-n}=\left(\dfrac{b}{a}\right)^{n}$ \ (si $a\neq0$). L'exponent negatiu no canvia el signe del resultat: \ $2^{-3}=\frac{1}{8}$, no $-8$.}

\ex{2} Ordena de menor a major: $\ (-2)^{-1},\ 2^{-2},\ (-2)^{-3},\ 2^{0},\ (-2)^{2},\ 2^{-1}$.

\ex{2} Expressa com una potència de base 2:\begin{tasks}[after-item-skip=8pt](4)
  \task $8^{-2}$
  \task $\left(\frac{1}{4}\right)^{3}$
  \task $\left(\frac{1}{16}\right)^{-2}$
  \task $\dfrac{4^{-1}\cdot2^{5}}{8^{2}}$
\end{tasks}

\ex{2} Expressa com una sola potència (amb exponents enters):\begin{tasks}[after-item-skip=8pt](3)
  \task $7^{-3}\cdot7^{5}$
  \task $3^{4}:3^{-2}$
  \task $(5^{-2})^{3}$
  \task $2^{-4}\cdot5^{-4}$
  \task $\left(\frac{3}{4}\right)^{-2}\cdot\left(\frac{3}{4}\right)^{5}$
  \task $6^{-3}:2^{-3}$
\end{tasks}

\ex{3} Redueix descomponent en factors primers i calcula:\begin{tasks}[after-item-skip=8pt](3)
  \task $\dfrac{9^{2}\cdot3^{-5}}{27^{-1}}$
  \task $\dfrac{4^{3}\cdot6^{-2}}{3^{-4}\cdot2^{5}}$
  \task $\dfrac{10^{-2}\cdot25^{2}}{5^{3}\cdot2^{-3}}$
\end{tasks}

\ex{3} Calcula pas a pas:\begin{tasks}[after-item-skip=10pt](2)
  \task $2^{-1}+2^{-2}-2^{-3}$
  \task $\left(\frac{2}{3}\right)^{-2}-\left(\frac{1}{2}\right)^{-1}\cdot \left(-3\right)^{-1}$
  \task $\left[\left(\frac{1}{2}\right)^{-2}-\left(-3\right)^{0}\right]^{-1}$
  \task $\left(\frac{3}{4}-\frac{1}{2}\right)^{-2}:\left(2-\frac{4}{3}\right)^{-1}$
\end{tasks}

\vfill
\begin{center}\color{muted}\small\faCheckCircle[regular]\ \ Ànim! Comprova els teus resultats amb les eines interactives dels apunts de la unitat.\end{center}

\end{document}
